\begin{table}[!htbp]
\centering\footnotesize\setlength{\tabcolsep}{4pt}
\caption{Чувствительность замороженной integer-KAN к точному совпадению тестовых строк с training по всем 44 исходным колонкам. Коэффициенты, LUT1025 и LUT513 дают одинаковые решения. Это постфактум анализ существующего теста без переобучения и изменения выбора.}
\label{tab:duplicate-sensitivity}
\begin{tabular}{lrrrrrr}
\toprule
Подгруппа & $N_0$ & $N_1$ & F1 & BA & FPR & TPR\\
\midrule
Весь тест & 10,000 & 32,209 & 0.982584 & 0.975391 & 0.021400 & 0.972182\\
Совпадающие & 1,634 & 3,567 & 0.996508 & 0.992350 & 0.015300 & 1.000000\\
Без точного совпадения & 8,366 & 28,642 & 0.980823 & 0.973063 & 0.022591 & 0.968717\\
\bottomrule
\end{tabular}
\end{table}

Качество на 5,201 совпадающей строке выше, чем на 37,008 остальных (табл.~\ref{tab:duplicate-sensitivity}). Различие не является причинной оценкой оптимизма: состав подгрупп различается, а обучение не повторялось. «Без точного совпадения» означает лишь отсутствие равной строки по всем 44 колонкам, включая метки и идентификаторы; совпадение только признаков отдельно не проверено. Подгруппа не является независимым тестом по узлам, сеансам или времени. Все три несертифицированных решения LUT1025 и все двенадцать LUT513 находятся в этой подгруппе; фактических различий с коэффициентным эталоном нет. Приложение~\ref{app:subgroup-details} сохраняет числа исходов и сертификации. Полный тест оценён на хосте по подготовленным входам; стоимость MCU измерена по отдельным протоколам и не объединяется с ним в единый тест точности, времени и энергии.
